\documentclass[preprint,12pt,a4paper]{elsarticle}
\usepackage{lmodern}
\usepackage{amssymb}
\usepackage{booktabs}
\usepackage{tabularx}
\usepackage{graphicx}
\usepackage{xcolor}
\usepackage{tikz}
\usepackage{float}
\usepackage{tabularx}
\usepackage{array}
\usepackage[table]{xcolor}
\usepackage{placeins}
\usepackage{hyperref}

\usetikzlibrary{arrows.meta,positioning,fit,shapes.geometric}
\setlength{\parindent}{0pt}
\setlength{\emergencystretch}{2em}
\newcolumntype{L}[1]{>{\raggedright\arraybackslash}p{#1}}
\newcolumntype{Y}{>{\raggedright\arraybackslash}X}
\newcommand{\tablerowsep}{\specialrule{0.25pt}{0.6pt}{0.6pt}}
\newcommand{\figurepanellabel}[1]{\par{\footnotesize #1}}
\makeatletter
\setlength{\@fptop}{0pt}
\setlength{\@fpsep}{10pt}
\setlength{\@fpbot}{0pt plus 1fil}
\setlength{\@dblfptop}{0pt}
\setlength{\@dblfpsep}{10pt}
\setlength{\@dblfpbot}{0pt plus 1fil}
\makeatother
\journal{SoftwareX}

% Generated by the release pipeline. The fallbacks keep draft builds readable,
% while check_softwarex_release.py rejects them for publication.
\IfFileExists{generated/release_metadata.tex}{\input{generated/release_metadata.tex}}{
  \newcommand{\ReleaseVersion}{v1.0.0 release candidate}
  \newcommand{\ReleaseDOI}{release gate pending}
  \newcommand{\ReleaseCommit}{release gate pending}
  \newcommand{\ReleaseDate}{release gate pending}
}

\begin{document}
\begin{frontmatter}

\title{Medora: A Consent-Aware, AI-Native Bilingual Platform for Healthcare Management and Consultation}

\author{Sarwad Hasan Siddiqui}
\author{Adiba Tahsin}
\author{Kazi Saeed Alam\corref{cor1}}
\ead{saeed.alam@cse.kuet.ac.bd}
\cortext[cor1]{Corresponding author.}
\author{Sk. Imran Hossain}
\address{Department of Computer Science and Engineering, Khulna University of Engineering \& Technology, Khulna 9203, Bangladesh}

\begin{abstract}
Medora is a bilingual (English/Bengali), consent-aware platform that combines patient-held records, appointment coordination, and 21 assistive-AI endpoints within an installable progressive web application. Patients can describe a health concern by text or voice to discover relevant doctors without needing to know a specific doctor’s name. Clinicians receive structured intake information, speech-to-notes support, source-linked patient histories, clinical-query assistance, and reviewable prescription drafts. Chorui provides role-aware chat, voice interaction, navigation, and Vapi-based live conversations across registered workflows. External processing requires revocable, purpose-specific consent and applies bilingual redaction and pseudonymization, while a fixed registry constrains assistant navigation. Transactional booking manages concurrent requests and post-commit notifications. Medora also provides an open Bangladesh medicine reference containing 7,389 canonical drugs, 67,001 brands, and 74,390 searchable terms. A prescription-processing prototype integrates a trained YOLO-based region detector with OCR, dosage parsing, and catalog matching. The software, data-building pipeline, tests, deployment configuration, and reproducibility artifacts are released as open research code. Medora is research software and is not a validated medical device.
\end{abstract}

\begin{keyword}
digital health \sep assistive AI \sep consent-gated processing \sep
medicine reference data \sep deterministic intent routing \sep bilingual software
\end{keyword}
\end{frontmatter}

\section*{Code metadata}
The code metadata required for reproducibility is summarized in [\hyperref[codeMetadata]{Table~\ref*{codeMetadata}}].

\begin{table}[!tb]
\small
\begin{tabular}{@{}L{0.05\textwidth}L{0.35\textwidth}L{0.53\textwidth}@{}}
\toprule
\textbf{Nr.} & \textbf{Description} & \textbf{Value} \\
\midrule
C1 & Current code version & \ReleaseVersion \\
\tablerowsep
C2 & Permanent link to code/repository used for this code version & \url{https://github.com/CSE-3200-System-Project/Medora/tree/\ReleaseCommit}; archive \ReleaseDOI \\
\tablerowsep
C3 & Permanent link to Reproducible Capsule & Not applicable; the archived repository contains the pinned containers, scripts, raw outputs, and verification receipt. \\
\tablerowsep
C4 & Legal Code License & MIT License \\
\tablerowsep
C5 & Code versioning system used & Git \\
\tablerowsep
C6 & Software code languages, tools, and services used & Python, TypeScript, SQL; FastAPI, Next.js, PostgreSQL; hosted text models (Groq, Gemini, Cerebras) with a deterministic mock; optional faster-whisper, Vapi, PaddleOCR, and Azure Document Intelligence \\
\tablerowsep
C7 & Compilation requirements, operating environments \& dependencies & Python 3.11, Node.js 20, PostgreSQL 16, or checksummed containers; reference deployment on Vercel and Azure Container Apps \\
\tablerowsep
C8 & If available Link to developer documentation/manual & \url{https://github.com/CSE-3200-System-Project/Medora/tree/main/docs} \\
\tablerowsep
C9 & Support email for questions & saeed.alam@cse.kuet.ac.bd \\
\bottomrule
\end{tabular}
\caption{Code metadata. The archived release date is \ReleaseDate.}
\label{codeMetadata}
\end{table}


\section{Motivation and significance}

Medora studies how bilingual language-model assistance can support healthcare workflows
while preserving review, consent, and clear system boundaries. Records, doctor discovery,
booking, summaries, voice, and medicine lookup share one authorization and audit model. The
platform targets health-informatics researchers, educators, and teams studying consent-aware
AI and reliable bilingual workflows.


Studies of Bangladesh's long-running digital-health program identify opportunities to
strengthen eHealth implementation and electronic-record adoption
~\cite{ahmed2014,khan2012}. Medora contributes an inspectable, mobile-first bilingual
artifact that connects patient-held information to search, appointments, and record views
while keeping the user in control of each action.

Three contributions provide the main reuse value. First, Medora supplies an open Bangladesh
medicine reference, addressing limited programmatically searchable resources. RxNorm, for
example, does not cover Bangladesh~\cite{rxnorm}. A reproducible build reconciles five
published datasets into generic, brand, and typo-tolerant search tables used throughout the
platform. The shared identity layer supports medicine lookup, prescription composition,
medication-history organization, and OCR matching without separate dictionaries for each
feature. The normalized corpus and transformation script make the curation reusable.


Second, patients can describe a condition by text or voice to obtain specialty-oriented
doctor candidates. Source-linked summaries organize patient and prescription histories,
while Chorui reaches registered workflows through chat, speech, or Vapi conversation. A
user can ask in ordinary language, clarify missing information, and continue into the
corresponding screen. Clinicians receive parallel intake, consultation, dictation, record-
review, and prescribing support.

Third, Medora provides a reusable containment pattern. Consent precedes external calls,
assistant navigation resolves through a fixed registry, and backend-attached references
limit summary citations to supplied records. Most safety-critical logic is deterministic
(Table~\ref{tab:components}). These controls allow hosted or local models to be evaluated
without giving them authority over access, routes, or clinical writeback.


OpenMRS, Bahmni, and GNU Health provide broad institutional record, hospital, laboratory,
and FHIR capabilities~\cite{openmrs,bahmni,gnuhealth}. Medora complements them with a
compact mobile client connecting patient-held records and transactional booking to
bilingual assistive AI (Table~\ref{tab:related-work}).

Research on summarization and symptom checkers motivates traceability and careful urgency
handling~\cite{shing2021,tracsum,wallace2022}. Medora links summaries to backend records,
presents specialty candidates rather than diagnoses, retains manual browsing, and applies
deterministic emergency guidance before model calls.


\begin{table}[!tb]
\scriptsize
\setlength{\tabcolsep}{3pt}
\begin{tabularx}{\textwidth}{@{}L{0.19\textwidth}YYYY@{}}
\toprule
\textbf{Criterion} & \textbf{OpenMRS}\cite{openmrs} & \textbf{Bahmni}\cite{bahmni} & \textbf{GNU Health}\cite{gnuhealth} & \textbf{Medora} \\
\midrule
Deployment & Modular server platform & Multi-service hospital distribution & Integrated hospital system & Three-service research stack \\
\tablerowsep
Patient mobile workflow & Deployment-specific web clients & Hospital-oriented web workflows & Web and desktop clients & English/Bengali patient PWA \\
\tablerowsep
Booking and live slots & Scheduling modules; deployment-specific updates & Appointment workflows & Appointment management & Transactional booking; post-commit updates \\
\tablerowsep
Assistive AI & Outside this comparison scope & Outside this comparison scope & Outside this comparison scope & 21 endpoints; text/voice assistance; source-linked drafts \\
\tablerowsep
Medicine reference & Deployment-supplied dictionaries & Deployment-supplied dictionaries & Deployment-supplied dictionaries & 74{,}390-term Bangladesh index; build script \\
\tablerowsep
Consent and audit & Privilege and audit modules & OpenMRS roles and audit tools & Role and audit tools & Purpose-, provider-, and scope-specific grants \\
\tablerowsep
Bengali localization & Localization framework & Localization support & Localization framework & Bundled English and Bengali catalogs \\
\tablerowsep
Interoperability & REST/FHIR ecosystem & OpenMRS interoperability ecosystem & FHIR and health standards & JSON/REST contracts; FHIR mapping is future work \\
\tablerowsep
Release status & Open source; mature deployments & Open source; hospital deployments & Open source; international deployments & MIT research release; validation pathway documented \\
\bottomrule
\end{tabularx}
\caption{Scope comparison from project documentation. Cells summarize documented emphasis and do not imply that unlisted capabilities are absent.}
\label{tab:related-work}
\end{table}

\section{Software description}

\subsection{Software architecture}

Medora comprises a Next.js web client, a FastAPI core API, a separate FastAPI OCR service,
and PostgreSQL (Fig.~\ref{fig:trust}). Authentication identifies the actor; backend checks
of role, ownership, care relationship, appointment state, and consent authorize each
operation. Frontend visibility is only a usability aid.

The reference deployment is available at
\url{https://medorahealth.vercel.app/}. Vercel serves the Next.js frontend, while the
FastAPI backend runs on Azure Container Apps. This public instance demonstrates the
software; the archived containers remain the reproducible release environment.

Release hardening covers every data path. An audit identified pre-release anonymous database
grants and disabled row-level security missed by API-only tests. The release revokes those
grants, enables row-level security, limits direct subscriptions to a trigger-maintained
change feed, and verifies these invariants with negative tests. Thus, authorization is
tested across both application and database interfaces.

\begin{figure}[!tb]
\centering\input{figures-src/trust_boundary.tex}
\caption{Deployment and trust boundaries. Local speech recognition and local OCR remain
inside the Medora boundary. Hosted text models, Vapi live audio, and Azure OCR are separate
recipients with separate grants.}
\label{fig:trust}
\end{figure}

Consent-aware processing is deny-by-default (Fig.~\ref{fig:consent}). Grants record subject,
recipient, purpose, scopes, policy version, validity, revocation, and audit metadata. Invalid
grants receive an explanation and, where available, a local alternative. Revocation blocks
future processing, and local mode never silently escalates to cloud processing. Provider-
and purpose-specific grants distinguish hosted text, Vapi audio, and Azure OCR processing.

\begin{figure}[!tb]
\centering\input{figures-src/consent_flow.tex}
\caption{Consent-aware processing. Purpose and provider are checked before any external
call, and outputs return as review-ready drafts.}
\label{fig:consent}
\end{figure}

\subsection{Software functionalities}

\textbf{Bangladesh medicine reference.} A reproducible script cleans five datasets into
three shared tables (Table~\ref{tab:corpus}). \texttt{drugs} stores identities keyed by
generic name, strength, and form; \texttt{brands} maps commercial names to them; and
\texttt{medicine\_search\_index} provides trigram-based typo tolerance. Search, prescribing,
histories, and OCR matching therefore use one identity. The build normalizes spelling,
strength, dosage-form, manufacturer, and duplicate-brand variations while preserving source
provenance. This produces a practical lookup layer and a repeatable basis for future source
refreshes. Dosage recommendations, interaction checking, and effectiveness ranking remain
separate future extensions.

\begin{table}[!tb]
\centering
\small

\renewcommand{\arraystretch}{1.15}
\setlength{\tabcolsep}{6pt}

\begin{tabular}{@{}lr@{}}

\toprule
\textbf{Artifact} & \textbf{Count} \\
\midrule

Consolidated source rows
& 71{,}795 \\[2pt]

Canonical drugs (generic $\times$ strength $\times$ form)
& 7{,}389 \\[2pt]

Brand entries
& 67{,}001 \\[2pt]

Search-index terms
& 74{,}390 \\[2pt]

Distinct generic names
& 5{,}242 \\[2pt]

Distinct brand names
& 52{,}117 \\

\bottomrule

\end{tabular}

\caption{Medicine reference as loaded. Counts differ from the consolidated CSV because the canonical key collapses brand duplicates.}

\label{tab:corpus}
\end{table}

\textbf{Assistive-AI layer.} Twenty-one endpoints cover doctor discovery, summaries, intake,
consultation, references, speech-to-notes, safety checks, and live voice
(Table~\ref{tab:ai-surface}). One provider-agnostic orchestrator supports Groq, Gemini,
Cerebras, and a deterministic mock. It sanitizes input, assigns a random correlation ID,
schema-validates responses, and records operational metadata, allowing providers to change
without rewriting workflows. Typed response models turn generated output into predictable
application contracts; malformed responses fail explicitly instead of entering forms or
records as unchecked prose.

Before orchestration, consent filters permitted categories and marks partial records; a
bilingual redactor removes known identifiers; and pseudonymization replaces patient IDs.
Raw clinical content is excluded from default logs. Clinician access also requires a care
relationship and appears in the patient's audit view. These stages make privacy and
authorization behavior observable in tests and reusable across all endpoints instead of
being reimplemented feature by feature.

\begin{table}[h]
\scriptsize
\setlength{\tabcolsep}{4pt}
\begin{tabularx}{\textwidth}{@{}>{\hsize=.80\hsize\raggedright\arraybackslash}X>{\hsize=.66\hsize\raggedright\arraybackslash}X>{\hsize=1.94\hsize\raggedright\arraybackslash}X>{\hsize=.60\hsize\raggedright\arraybackslash}X@{}}
\toprule
\textbf{Component} & \textbf{Kind} & \textbf{Role} & \textbf{Model} \\
\midrule
AI orchestrator & Interface & Provider-agnostic call path: sanitize, dispatch, schema-validate, log latency & Yes \\
\tablerowsep
Consent guard & Deterministic & Resolves permitted data categories; filters payload; emits restriction notice & No \\
\tablerowsep
Bilingual redactor & Deterministic & Pattern removal of contact, identity, address, and date spans in Bengali and English & No \\
\tablerowsep
Patient pseudonymiser & Deterministic & Stable opaque reference replacing raw identifiers & No \\
\tablerowsep
Processing-consent service & Deterministic & Purpose-, provider-, and scope-scoped external grants with revocation & No \\
\tablerowsep
Chorui intent normalizer & Deterministic & Utterance to canonical intent, bilingual & No \\
\tablerowsep
Chorui route registry & Deterministic & Immutable role-scoped intent-to-route table; admin routes absent & No \\
\tablerowsep
Chorui navigation engine & Deterministic & Role-aware resolution, missing-parameter clarification, fallbacks & No \\
\tablerowsep
Specialty matcher & Deterministic & Exact, synonym, fuzzy, and similarity matching of specialty terms & No \\
\tablerowsep
Medical knowledge base & Deterministic & Symptom-to-specialty relations and universal fallback chains & No \\
\tablerowsep
Emergency red-flag rules & Deterministic & Keyword screen that pre-empts model output with emergency guidance & No \\
\tablerowsep
Speech recognition (faster-whisper) & Local model & Bilingual transcription of clinician dictation & Yes, local \\
\tablerowsep
Live-voice bridge (Vapi) & External & Tool-calling surface resolving through the same registry and guards & Yes, external \\
\tablerowsep
Dosage matcher & Deterministic & Dosage, frequency, duration, and quantity grammar & No \\
\tablerowsep
Medicine fuzzy matcher & Deterministic & Trigram and edit-distance matching against the reference index & No \\
\tablerowsep
Region detector (YOLO) & Local model & Prescription region localisation & Yes, local \\
\tablerowsep
Local OCR (PaddleOCR) & Local model & On-premise text recognition & Yes, local \\
\tablerowsep
Cloud OCR (Azure) & External & Optional recognition backend under a separate grant & Yes, external \\
\tablerowsep
Report and prescription parsers & Deterministic & Structured extraction of test results and medication rows & No \\
\bottomrule
\end{tabularx}
\caption{AI and machine-learning components. Thirteen of nineteen use deterministic logic;
the remaining six coordinate or perform language, speech, and image-model inference.}
\label{tab:components}
\end{table}

Figure~\ref{fig:chorui} summarizes the assistant's deterministic request path.

\begin{figure}[!tb]
\centering\includegraphics[width=\textwidth]{figures-src/chorui_architecture.pdf}
\caption{Chorui request path. Text or voice requests receive conversational phrasing while
deterministic intent normalization and registry resolution select role-appropriate app
destinations.}
\label{fig:chorui}
\end{figure}

\begin{table}[!t]
\centering
\scriptsize

\renewcommand{\arraystretch}{1.18}
\setlength{\tabcolsep}{4pt}

\begin{tabularx}{\textwidth}{
@{}
>{\raggedright\arraybackslash}p{0.23\textwidth}
>{\centering\arraybackslash}p{0.04\textwidth}
>{\raggedright\arraybackslash}X
>{\raggedright\arraybackslash}p{0.15\textwidth}
@{}
}

\toprule
\textbf{Endpoint group} &
\textbf{No.} &
\textbf{User-facing contribution} &
\textbf{Verification} \\
\midrule

Assistant chat and conversations
& 4
& Bilingual help, conversation history, and guided navigation
& Navigation fixtures \\[2pt]

Specialty and doctor search
& 1
& Condition-to-specialty doctor candidates
& Navigation fixtures \\[2pt]

Record summarisation
& 3
& Source-linked patient and prescription histories
& Summary fixtures \\[2pt]

Intake structuring
& 2
& Confirmable structured intake drafts
& Contract + auth. \\[2pt]

Consultation assistance
& 6
& Reviewable notes, queries, summaries, and prescription drafts
& Contract + auth. \\[2pt]

Clinical reference query
& 1
& Structured reference text with confidence and cautions
& Contract + auth. \\[2pt]

Speech to notes
& 1
& Bilingual dictation to editable clinical forms
& Contract + auth. \\[2pt]

Safety check
& 1
& Advisory flags inside the clinician workflow
& Contract + auth. \\[2pt]

Live voice tool calls
& 2
& Vapi conversation with role-aware tools
& Contract + auth. \\[3pt]

\midrule

\textbf{Total}
& \textbf{21}
&
& \textbf{All checked; 3 fixture groups} \\

\bottomrule

\end{tabularx}

\caption{The assistive-AI surface. Every group shares consent, redaction,
pseudonymisation, schema validation, contract checks, and authorization tests;
navigation and summary groups additionally have task-specific fixtures.
Full paths, schemas, authorization rules, and example payloads are retained
in the archived API specification.}

\label{tab:ai-surface}
\end{table}

\FloatBarrier

\textbf{Chorui assistant.} Chorui accepts Bengali or English text and voice requests and
guides users to role-appropriate workflows. An immutable registry binds 23 canonical
intents to routes, required parameters, and fallbacks through 24 entries: 11 patient, 11
clinician, and 2 shared. Missing details trigger clarification, while unregistered
administrative routes remain unreachable. Reachable actions are fixed at build time and
auditable in one file. The same resolution path serves chat and Vapi tool calls.

\textbf{Summaries and speech.} Medora organizes patient, clinician, and prescription
histories into role-specific, read-only summaries. Backend-attached references keep every
item traceable to supplied records and allow return to the underlying entry. Empty source
sets produce a not-found result, and summaries never write back automatically. Local
faster-whisper supports bilingual dictation~\cite{whisper}; Vapi adds live conversation and
the same Chorui and doctor-search tools under a separate live-audio grant.

All 21 endpoints receive contract and authorization coverage; navigation and summary groups
also have reproducible behavioral fixtures. This separates interface assurance from
task-specific measurement and supports further endpoint evaluation.

\textbf{Human-centered workflow.} Summaries remain reviewable, prescription suggestions
pre-fill fields a prescriber edits and signs, and specialty search retains direct browsing.
OCR produces row-by-row drafts marked \texttt{authoritative\_writeback=false}; patients may
acknowledge receipt or report discrepancies without implying clinical endorsement. These
review points accelerate organization while leaving decisions and changes with users.

\textbf{Installable bilingual client.} The PWA supports English and Bengali, standalone
display, maskable icons, and home-screen installation. Browser tests cover both locales.
Only public/versioned resources are cached; health writes stay online, and logout clears
caches, IndexedDB, and sensitive storage.

\textbf{Reliable appointments.} Idempotency keys, row locks, availability checks, and a
uniqueness constraint support concurrent booking. Appointment, audit, replay result, and
outbox event commit together; clients then receive updates and recover PostgreSQL state
after reconnection. Reminders use in-app notifications, email, and optional web push. This
connects condition-driven discovery to a transactionally protected appointment workflow
rather than ending at an AI recommendation screen.

\IfFileExists{generated/safety_results.tex}{\input{generated/safety_results.tex}}{}

\subsection{Evaluation and reproducibility}

Evaluation combines privacy, navigation, summary, and booking benchmarks with contract,
authorization, unit, integration, browser, and container checks. Mock and live results are
reported separately. The tables are generated from archived machine-readable outputs, so
their values can be regenerated rather than manually transcribed into the manuscript. Tables~\ref{tab:safety-results}--\ref{tab:booking-results}
present the evaluation results, covering overall safety fixture outcomes,
symptom-navigation performance, span-level privacy metrics, and booking
contention under concurrent requests.

The synthetic bilingual privacy corpus covers identifiers, spelling variants, obfuscation,
over-redaction, consent, and prompt injection (Table~\ref{tab:privacy-span-results}). Written
independently of the patterns, it measures 94.7\% precision and 75.5\% recall. Every miss is
documented, exposing unlabelled names, month-name dates, and obfuscated contacts as concrete
improvement targets.

Navigation fixtures exercise recorded-provider and deterministic-mock paths, verifying
manual browsing, catalog filtering, and conservative emergency responses. Summary fixtures
verify backend-supplied citations, empty-input behavior, and rejection of malformed output.

Booking tests run 2, 10, and 50 concurrent attempts over 30 fresh slots per level. Each run
requires one committed appointment, no duplicate, idempotent replay, rollback-safe retry,
and post-commit outbox processing.

The provider manifest records model and interface settings, region, retries, retention, and
paid-service requirements. The archive includes commands, checksums, raw outputs, and table
generation instructions. A deterministic provider reproduces schemas, refusals, and failure
paths without network credentials, while live-provider results remain separately identified.
It supports reproducible comparison across provider configurations.
%

\IfFileExists{generated/booking_results.tex}{\input{generated/booking_results.tex}}{}

\textbf{Prescription image prototype.} A trained YOLO model isolates prescription regions
for PaddleOCR or Azure Document Intelligence~\cite{du2020,prescriptionocr}. Bengali-numeral
normalization, dosage grammars, structured parsing, and fuzzy matching against 67,001 brands
complete the image-to-catalog workflow. Each extracted row requires confirmation before a
record update.

Cursive handwriting does not yet support an accuracy claim: partial names may yield an
incorrect match or no match at stricter thresholds. The release preserves the models,
services, parsers, prelabels, and review workflow for a future clinically adjudicated
benchmark; review gating keeps prototype output separate from medication records. The
completed contribution is therefore the reproducible region-to-catalog pipeline and its
integration with the curated reference, with handwriting recognition identified as the
next measurable research stage.

\section{Illustrative examples}

Figures~\ref{fig:ui-bilingual}--\ref{fig:ui-ai} show the demonstration interfaces used by
these fixtures; their provenance is described in the Ethics section.

\begin{figure}[!tb]
\centering
\includegraphics[width=0.92\textwidth]{figures-ui/ui_dashboard_en.png}\\[3pt]
\figurepanellabel{(a) English}\vspace{6pt}
\includegraphics[width=0.92\textwidth]{figures-ui/ui_dashboard_bn.png}\\[3pt]
\figurepanellabel{(b) Bengali}
\caption{A consistent record view across the two bundled locales. The shared layout and
localized message catalog provide the same patient workflow in English and Bengali;
additional server-generated strings can be added to the catalog incrementally.}
\label{fig:ui-bilingual}
\end{figure}

\begin{figure}[!tb]
\centering
\includegraphics[width=0.80\textwidth]{figures-ui/ui_consent_sharing.png}\\[3pt]
\figurepanellabel{(a) Per-recipient, per-category sharing}\vspace{6pt}
\includegraphics[height=6.0cm]{figures-ui/ui_consent_audit.png}\hfill
\includegraphics[height=6.0cm]{figures-ui/ui_mobile_consent.png}\\[3pt]
\figurepanellabel{\makebox[\textwidth]{(b) Patient-visible access log\hfill(c) Sharing controls at a phone viewport}}
\caption{The consent surface. Patients grant categories per recipient, revise them at any
time, and review clinician or assistant-mediated access in the audit log. Panel (c) shows
the same controls in the installed mobile client.}
\label{fig:ui-consent}
\end{figure}

\begin{figure}[!tb]
\centering
\includegraphics[width=0.80\textwidth]{figures-ui/ui_ai_navigation.png}\\[3pt]
\figurepanellabel{(a) Patient-facing specialty navigation}\vspace{6pt}
\includegraphics[height=6.0cm]{figures-ui/ui_ai_summary.png}\hfill
\includegraphics[height=6.0cm]{figures-ui/ui_mobile_assistant.png}\\[3pt]
\figurepanellabel{\makebox[\textwidth]{(b) Clinician-facing record summary\hfill(c) Assistant at a phone viewport}}
\caption{Assistive AI across patient and clinician roles. A typed or spoken condition can
launch specialty-oriented doctor discovery in (a); (b) organizes supplied records into a
source-linked clinician summary; and (c) provides mobile conversational access to Chorui.}
\label{fig:ui-ai}
\end{figure}

\textbf{Condition-to-doctor discovery.} A patient can type or speak a current concern in
the assistant panel without knowing a doctor's name (Fig.~\ref{fig:ui-ai}a). For the worked
example, ``Patient name: Test Rahim; phone: 01700000000; I have headache and fever'' becomes
``Patient name: [redacted-name]; phone: [redacted-phone]; I have headache and fever'' before
the schema-gated request. Medora converts the concern into specialty-oriented candidates,
opens the corresponding doctor search, communicates uncertainty, and keeps manual browsing
available. The same interaction can be entered through local speech recognition or a
consented Vapi conversation. Inputs, authorization state, schema, and outputs are archived
in \texttt{worked\_examples.json}.

\textbf{Registry-guided assistance.} Chorui resolves everyday requests through the same
role-scoped registry used by text and voice tools. It can clarify missing details, preserve
conversation history, and guide users to registered patient or clinician workflows. A
boundary test asks it to open the administrator dashboard; because that destination is not
registered for the patient role, Chorui safely returns its fallback even when prompt
injection names the route directly.

\textbf{Consent-scoped summarization.} A clinician requests an organized summary for a
patient who shared conditions and medications but not test results
(Fig.~\ref{fig:ui-consent}a). The guard includes the permitted categories, marks the
laboratory category as unavailable, and produces a source-linked summary that reports the
scope clearly. The patient can see the access in the audit table.

\section{Impact}

Medora enables research questions that are difficult to study with isolated demonstrations:
how role-bounded assistants change navigation effort; how consent granularity affects useful
context; how bilingual redaction behaves across scripts; and how source attachment,
transaction boundaries, and human review interact in one application. Its deterministic
registry and mock provider allow these questions to be examined while holding route access
and provider behavior stable.

For daily workflows, patients gain bilingual record management, condition-driven doctor
discovery, transactional booking, reminders, consent controls, and text or voice assistance.
Clinicians gain organized patient and prescription summaries, speech-to-notes, structured
intake, consultation drafts, clinical queries, prescription support, and visible correction
paths. Vapi extends these capabilities into live conversation, while the PWA makes them
installable on mobile devices. The public reference deployment lets researchers and
educators inspect the integrated experience without first assembling the service stack.
It also supports demonstrations, reviewer inspection, and reproducibility checks.

The reusable contributions extend beyond the interface. The Bangladesh medicine reference
and its build script can support independent medicine-search projects; the trained YOLO,
OCR, parser, and matcher pipeline offers a reproducible foundation for prescription-image
research; the deterministic mock enables tests without paid credentials; and the
role-scoped registry, consent state machine, source attachment, booking transaction, and
outbox patterns transfer to other data-sensitive applications. Because these assets share
documented contracts but remain modular, researchers can reuse the medicine corpus, privacy
pipeline, assistant registry, or concurrency design without adopting the complete product.
The release makes no unsupported claim of widespread or commercial adoption; its immediate
impact is a public, citable foundation for replication, teaching, comparative evaluation,
and collaboration.

\textbf{Current scope and next steps.} The release provides task-specific fixtures for
navigation and summaries and contract/authorization coverage for the remaining endpoint
groups; future work can add domain-specific accuracy studies for intake, consultation,
reference, speech, safety, and live voice. The medicine reference is a major consolidation
of the available Bangladesh-oriented sources and is published with provenance and a
reproducible build; future collaboration with regulatory and clinical experts can further
review indication text and refresh coverage. The OCR subsystem already establishes the
complete region-to-catalog workflow, with clinically adjudicated handwriting evaluation as
its next milestone. Redaction rules, emergency patterns, and provider assumptions are
explicitly versioned so they can continue to improve. The reported concurrency results
describe the tested environment. Clinical trials, regulatory assessment, and penetration
testing remain necessary before clinical deployment.
These boundaries do not reduce the software contribution; they define an evidence-based
path from an open research platform to future clinical evaluation.

\section{Conclusions}
Medora integrates bilingual patient-held records, condition-driven doctor discovery, transactional appointment coordination, source-linked histories, clinical drafting, Chorui text and voice assistance, Vapi-based live conversations, and prescription-image research within a single inspectable platform. Its principal reusable assets include the open Bangladesh medicine reference and data-build pipeline, trained region-detection and OCR integration, and the consent, navigation, traceability, and concurrency mechanisms that constrain assistive AI by design. Together, these components provide a robust foundation for software research, reproducibility, and future clinical evaluation. 

\section*{Ethics statement and data availability}
\label{sec:ethics}

The Medora source code is released under the MIT License. Annotations, synthetic
test fixtures, and aggregate evaluation results are released under CC BY 4.0.
Interface figures were generated from a non-production demonstration environment.
The patient-facing record shown in these figures belongs to a consenting author,
and all clinician profiles are test accounts.

The Bangladesh medicine reference consolidates five redistributable sources into
a normalized, search-oriented corpus. The source datasets are distributed under
CC0 1.0 (two sources), Apache 2.0, MIT, and CC BY 4.0 licenses~\cite{bdmeds}.
The consolidated Medora reference is released under CC BY 4.0 while retaining
source-specific attribution, provenance information, transformation logic, and
the notices required by the original licenses, including the Apache 2.0 statement
of changes. Two source datasets were originally derived from public websites;
therefore, the resulting corpus is presented as a transparent research
consolidation rather than as an official medicine register.

Prescription images used during development are not included in the public
release. These images contain identifiable clinical information and were
collected with verbal consent for an undergraduate research project, without
institutional ethics approval for public redistribution. Because names,
registration numbers, clinic information, and potentially identifying handwriting
may remain recoverable even after conventional redaction, the images are excluded
from the archived research artifacts. Any future release would require documented
consent and, where appropriate, institutional ethics approval or a controlled-access
data-sharing process.

A licensed clinician independently reviewed all 30 symptom-navigation fixtures.
During this review, two Bengali red-flag presentations were reclassified as
emergencies. The corresponding deterministic rules were subsequently updated, and
the complete review and revision trail is retained with the reproducibility
artifacts.

\section*{CRediT authorship contribution statement}
\textbf{Sarwad Hasan Siddiqui:} Conceptualization,  Software, Data curation, Methodology, Validation, Visualization, Writing -- original draft. \textbf{Adiba Tahsin:} Software, Methodology, Data curation, Writing -- original draft. \textbf{Kazi Saeed Alam:} Conceptualization, Supervision, Resources, Validation, Writing -- review and editing. \textbf{Sk. Imran Hossain:} Supervision, Writing -- review and editing.

\section*{Declaration of Generative AI and AI-assisted technologies in the writing process}
During the preparation of this work, the authors used ChatGPT in order to improve the readability and language of the manuscript. After using this service, the authors reviewed and edited the content as needed and take full responsibility for the content of the published article.

\section*{Declaration of competing interest}
The authors declare that they have no known competing financial interests or personal relationships that could have appeared to influence the work reported in this paper.

\section*{Funding}
This research received no specific grant from funding agencies in the public,
commercial, or not-for-profit sectors, and was self-funded.

\begin{thebibliography}{00}
\bibitem{openmrs} OpenMRS Community, Architecture, \url{https://devmanual.openmrs.org/technology/architecture/}, accessed 2026.
\bibitem{bahmni} Bahmni Coalition, Bahmni: Open Source EMR \& Hospital Information System, \url{https://www.bahmni.org/}, accessed 2026.
\bibitem{gnuhealth} GNU Solidario, GNU Health Hospital Information System documentation, \url{https://docs.gnuhealth.org/his/}, accessed 2026.
\bibitem{ahmed2014} T. Ahmed et al., eHealth and mHealth initiatives in Bangladesh: a scoping study, BMC Health Services Research 14 (2014) 260, \href{https://doi.org/10.1186/1472-6963-14-260}{doi:10.1186/1472-6963-14-260}.
\bibitem{khan2012} S.Z. Khan et al., Hopes and fears in implementation of electronic health records in Bangladesh, Electronic Journal of Information Systems in Developing Countries 54 (2012) 1--18, \href{https://doi.org/10.1002/j.1681-4835.2012.tb00387.x}{doi:10.1002/j.1681-4835.2012.tb00387.x}.
\bibitem{rxnorm} S. Liu, W. Ma, R. Moore, V. Ganesan, S. Nelson, RxNorm: prescription for electronic drug information exchange, IT Professional 7 (2005) 17--23, \href{https://doi.org/10.1109/MITP.2005.122}{doi:10.1109/MITP.2005.122}.
\bibitem{bdmeds} M.M. Rahman, M.M. Khan, Medicinal products in Bangladesh: a dataset of generic and brand names, dosage strengths, and manufacturers, Mendeley Data V1 (2024), \href{https://doi.org/10.17632/zhtvkny53n.1}{doi:10.17632/zhtvkny53n.1}.
\bibitem{whisper} A. Radford, J.W. Kim, T. Xu, G. Brockman, C. McLeavey, I. Sutskever, Robust speech recognition via large-scale weak supervision, in: Proc. ICML, 2023, pp. 28492--28518.
\bibitem{du2020} Y. Du et al., PP-OCR: A practical ultra lightweight OCR system, arXiv:2009.09941 (2020).
\bibitem{prescriptionocr} A. Datta et al., Preserving medical information from doctor's prescription ensuring relation among the terminology, Computers in Biology and Medicine 188 (2025) 109812, \href{https://doi.org/10.1016/j.compbiomed.2025.109812}{doi:10.1016/j.compbiomed.2025.109812}.
\bibitem{shing2021} H.-C. Shing et al., Towards clinical encounter summarization: Learning to compose discharge summaries from prior notes, arXiv:2104.13498 (2021).
\bibitem{tracsum} B. Chu, M. Li, S. Frihat, C. Gu, G. Lodde, E. Livingstone, N. Fuhr, TracSum: a new benchmark for aspect-based summarization with sentence-level traceability in medical domain, in: Proc. EMNLP, 2025, pp. 844--864, \href{https://doi.org/10.18653/v1/2025.emnlp-main.43}{doi:10.18653/v1/2025.emnlp-main.43}.
\bibitem{wallace2022} W. Wallace et al., The diagnostic and triage accuracy of digital and online symptom checker tools: a systematic review, npj Digital Medicine 5 (2022) 118.

\end{thebibliography}

\end{document}
